\begin{textsample}{POS Dim 1 – vem\_ed\_13 – Score 15.00 – vem\_ed\_13/t055.txt}  \label{ex:f1_pos_123}
EDUCAÇÃO INTERNACIONAL EM DEBATE Entre abril e junho de 2022, Osvaldo Succi Junior, coordenador dos Projetos Colaborativos Internacionais (PCIs / Cesu), divulgou experiências e \textbf{aprendizados} dos PCIs para a comunidade acadêmica.
Em 25 de abril, apresentou trabalho na Conferência 2022 da FAUBAI (Associação Brasileira de Educação Internacional).
Resumiu a história dos \textbf{intercâmbios} virtuais nas Fatecs, desde a primeira colaboração, em 2013: “Aproximando \textbf{futuros} gestores internacionais”, entre Fatec Americana e SUNY Ulster (EUA), venceu o Prêmio Santander Universidades.
Succi também expôs aspectos \textbf{que} contribuíram para a trajetória de sucesso dos PCIs nas Fatecs: um dos principais \textbf{fatores} externos \textbf{é} a pandemia do COVID-19, \textbf{que} aumentou o interesse por soluções tecnológicas para o ensino-aprendizagem (atualmente, a adaptação ao hibridismo presencial e virtual agrega novos desafios para projetos de \textbf{intercâmbio} virtual).
Entre os \textbf{fatores} internos destaca-se a defesa, por \textbf{parte} do Centro Paula Souza, em manter aulas de inglês e espanhol como \textbf{parte} dos currículos dos cursos.
Some-se a isso o \textbf{apoio} da Unidade do Ensino Superior de Graduação (Cesu) na \textbf{ação} junto às Regionais e às Unidades de Fatec e uma visão diferenciada na gestão dos PCIs, \textbf{que} \textbf{cria} laços perenes com parceiros \textbf{estabelecidos} e busca novas oportunidades nas instituições estrangeiras, \textbf{gerando} um fluxo constante de projetos.
Como resultado de persistência, \textbf{ações} planejadas e engajamento da equipe dos PCIs, a Unidade do Ensino Superior de Graduação (Cesu) do Centro Paula Souza alcançou, somente em 2021, 1.800 alunos de Fatecs envolvidos, com mais de 20 instituições de ensino \textbf{superior} internacionais, em 69 PCIs (mundialmente conhecidos como Virtual Exchange – \textbf{intercâmbios} virtuais – ou COIL – Collaborative Online International Learning).
Em 10 de maio, Succi Junior fez palestra na Fatec São Caetano do Sul, a convite da diretora da unidade, Adriane Monteiro Fontana.
Sintetizou a experiência da Cesu com Intercâmbios Virtuais, ressaltou a \textbf{importância} de sair da zona de conforto para lidar com os parceiros, “afiar” percepções culturais e \textbf{habilidades} \textbf{linguísticas}.
Em 2 e 3 de junho, participou por videochamada do “Taller Internacional de Docencia” organizado pela Uniminuto (Colômbia), cujo tema \textbf{foi} qualidade nos projetos COIL.
Presencialmente, na Colômbia, estavam Stephanie Doscher (Florida International University – FIU / EUA) e Javier Guerrero (Uniminuto).
Succi Junior ressaltou a \textbf{importância} das colaborações Sul-Sul: países diferentes, desafios semelhantes.

% matched lemmas: apoio, aprendizado, ação, criar, estabelecer, fator, futuro, gerar, habilidade, importância, intercâmbio, linguístico, parte, que, ser, superior
\end{textsample}
